\section{The Takeover in a Live Account}
\label{sec:case-study}

The corpus says the path is not completed in public infrastructure as code.
It does not say the path works. We built it in an AWS account to show that it
does, and to find out what actually gates it.

\subsection{Setup}

The account holds three agent runtimes and three execution roles. Two of the
roles differ in their policies. The third, \texttt{execC}, is identical to the
second in every way except its name.

The attacker is a separate principal whose only permission is the AWS managed
policy \iam{BedrockAgentCoreFullAccess}. It is given nothing else. In
particular it is given no direct access to the resource it will end up
reading.

A canary object sits in a bucket that only the victim's execution role can
read. A pre-check confirms the attacker cannot read it.

\subsection{Verdict rules}

Three rules are encoded in the harness. A verdict comes from reading the
runtime back from the platform and asking what role it says the runtime now
runs as, never from the request that was sent. An \texttt{AccessDenied} counts
as a block only when the message names \iam{iam:PassRole}; any other denial is
recorded as inconclusive, because a denial for an unrelated reason would
otherwise look like the control working. And role assumption is retried for
120 seconds, so a propagation delay is not recorded as a control.

\begin{figure}[t]
\centering
\begin{tikzpicture}[
  font=\small,
  atk/.style={draw, rounded corners=2pt, align=center, inner sep=3pt,
              text width=44mm, draw=cStored, fill=cStored!8, line width=0.7pt},
  same/.style={draw, rounded corners=2pt, align=center, inner sep=3pt,
               text width=26mm, minimum height=7mm, line width=0.5pt,
               draw=cIdentity, fill=cIdentity!8},
  verdict/.style={font=\scriptsize, align=center, inner sep=1pt,
                  text width=32mm},
  flow/.style={-{Stealth[length=1.8mm]}, line width=0.6pt},
]

\node[atk] (atk)
  {attacker principal\\[1pt]
   \scriptsize holds \iam{BedrockAgentCoreFullAccess}\\[-1pt]
   \scriptsize as published, nothing else};

% placed by explicit offset from the centre: anchoring to the box corners
% spreads them wider than the column
\node[same] (b) at ([shift={(-19mm,-13mm)}]atk.south) {\texttt{execB}};
\node[same] (c) at ([shift={( 19mm,-13mm)}]atk.south) {\texttt{execC}};

\draw[flow] (atk.south) -- (b.north);
\draw[flow] (atk.south) -- (c.north);

\node[verdict, below=1.5mm of b, text=cNeutral]
  {name does not match\\[1pt]\textbf{denied}};
\node[verdict, below=1.5mm of c, text=cStored]
  {name matches the pattern\\[1pt]\textbf{taken over}};

\node[font=\scriptsize, align=center, text width=80mm]
  at ([shift={(0mm,-32mm)}]atk.south)
  {Identical in policy, in trust and in permissions.\\
   The only difference is the string that names them.};

\end{tikzpicture}
\caption{The control. Two victim roles are built identically and differ only
in name. The attacker holds one AWS managed policy, attached as published,
with no inline statement added. One role is taken over and the other is
refused, so what separates one agent from another agent's role is a check on
the kind of role rather than on whose role it is, and the account's own
default role names satisfy that check.}
\label{fig:control}
\end{figure}


\subsection{The cells}

The experiment is nine cells, each a question with an outcome recorded to a
file at the time it ran. Table~\ref{tab:cells} lists them. The design is a
control and a treatment that differ by exactly one permission, plus the cells
that close off the ways the result could be an artefact.

\begin{table}[t]
\centering
\caption{The nine cells of the live experiment and the verdict each
recorded. The attacker holds the managed policy in every cell except the
first. Control and treatment differ by one inline grant and nothing else.}
\label{tab:cells}
\small
\begin{tabular}{p{18mm}p{39mm}p{15mm}}
\toprule
\textbf{Cell} & \textbf{Question} & \textbf{Verdict} \\
\midrule
positive control & can the operator replace the runtime's artifact at all &
  allowed \\
control & can the attacker, holding only the managed policy, replace a
  sibling's artifact & denied \\
treatment & the same, with one inline \iam{iam:\allowbreak PassRole} grant
  on the target role added & allowed \\
recheck & the same again with that grant removed & denied \\
\midrule
name condition & the same as the control, changing only the victim role's
  \emph{name} to match \iam{*Bedrock\allowbreak AgentCore*} & allowed \\
\midrule
gain & does the target role reach something the attacker does not & gain \\
effect & does the replaced code run and return the canary & canary
  returned \\
falsification & give the attacker directly what the target holds, then take
  over anyway & allowed, no gain \\
\midrule
persistence & after the attacker loses every permission, does the operator's
  own call still reach the planted code & persists \\
\bottomrule
\end{tabular}
\end{table}

The control and the recheck both end in an implicit deny, and the policy
simulator reports no matching statement. The treatment differs from them by a
single inline grant, and the recheck removes that grant and recovers the deny,
so the allow is attributable to the grant rather than to a propagation delay
or to a change in the account between runs.

The falsification cell is the one that constrains the predicate rather than
the platform. The attacker is given, directly, the thing the target role
holds. The takeover still succeeds and the API exchange is identical to the
treatment's. Equation~\ref{eq:esc} does not fire, because the reach is already
inside $\auth$.

\subsection{Result}

The attacker replaced the sibling runtime's definition. The platform reported
the runtime as running under the victim's execution role. The replaced code
executed, read the canary, and returned it to the attacker, who could not read
it directly before or after.

The planted code carries no SDK, because the runtime image has none, so it
signs its own request with a short standard library implementation. It
reports the \emph{name} of its credential source and the \emph{names} of the
AWS environment variables it can see, never a value.

\subsection{What the gate turned out to be}

The most useful cell is the name condition, because it changes one thing. The
victim role in that cell is identical to the control's in its policies, its
trust policy and its permissions. Only its name differs, and it differs so as
to match \iam{*BedrockAgentCore*}, which is how the AWS console names the
execution roles it creates for this service.

With that single change, and with the attacker still holding only
\iam{BedrockAgentCoreFullAccess} and no inline grant at all, the takeover
succeeded. The control denied the same request against the same kind of role
under a different name.

So the control that stands between two agents in a default deployment is a
string match on a role name, not a statement about which agent owns which
role. An operator who follows the console's own defaults, and grants the
vendor's own managed policy, has the path open.

This is a result about the condition rather than about the policy being
permissive. A policy named for full access that passed any role without a
condition would carry no claim to separate anything, and there would be
nothing here to report. The \iam{iam:PassRole} statement in v21 is not
unscoped: it is narrowed twice, by a resource pattern over the role's name
and by a condition on the service the role may be passed to, and narrowing a
statement is a decision to restrict what the permission reaches. What the
experiment establishes is that the restriction chosen, a glob over the name, is
not a restriction on ownership, and that the service's own console produces
names on the inside of it. The same statement with the role's ARN in place of
the pattern refuses the cell that succeeded, which
Section~\ref{sec:arn-repair} measures rather than asserts.

\subsection{What the second permission buys}

The canary could also have been read through the command route, with one
permission and no \iam{iam:PassRole}. The ninth cell measures what the extra
permission buys. It runs after the name condition, so the attacker holds only
the managed policy, and it aims the planted code at the canary through the
environment variables of the same \api{Update\allowbreak AgentRuntime} call.
Then every policy on the attacker principal is detached. The attacker's next
update of that runtime is refused with an \texttt{AccessDenied} naming
\iam{bedrock-agentcore:\allowbreak Update\allowbreak AgentRuntime}, and the
simulator reports an implicit deny on the canary.

With the attacker in that state, the operator invokes the runtime through its
ordinary endpoint. The response carries the attacker's marker, names instance
metadata as its credential source, and contains the canary value.

A stolen credential stops working when the principal that took it is disabled.
A replaced artifact keeps answering the agent's own callers under the victim's
role, with no attacker principal involved, until someone compares the running
artifact against the one that was meant to be there. That comparison is the
detection in Section~\ref{sec:defences}.
